\documentclass[12pt,a4paper,oneside]{article}
\usepackage[utf8]{inputenc}
\usepackage{amsmath}
\usepackage{tcolorbox}
\usepackage{amsfonts}
\usepackage{amssymb}
\usepackage{graphicx}
\usepackage[width=2.50cm, height=2.50cm, left=2.50cm, right=2.50cm, top=2.50cm, bottom=2.50cm]{geometry}
\author{HOUNTONDJI Sèmèvo Ronald Sylvanus, alias Gon FREECCSS}
\title{TP 1 : Structure d'un document LaTeX}
\begin{document}
\maketitle
\par Oh mon modèle, avec toi ma vie est toujours plus belle et auprès de toi, je me sens pousser des ailes.Te désobéir serait criminelle.  Vivre sans toi serait tellement cruelle car tu es quelqu'un d'exceptionnelle. \\
\rule{16cm}{1pt} \\ \\
{\Large \underline{\textbf{TP 2 :}} Mise en forme du texte} \\ 
\par \textbf{Oh bébé t'es mon bébé ! Toi qui m'a tout donné, je ne pourrai jamais te quitter je t'en fais la promesse mon bébé. Tu m'as tellement épaulé quand il le fallait. Tu m'as sorti de la galère je ne fais que rendre la monnaie.} \\
\par \textit{Oh bébé t'es mon bébé ! Toi qui m'a tout donné, je ne pourrai jamais te quitter je t'en fais la promesse mon bébé. Tu m'as tellement épaulé quand il le fallait. Tu m'as sorti de la galère je ne fais que rendre la monnaie.} \\ 
\par \underline{Oh bébé t'es mon bébé ! Toi qui m'a tout donné, je ne pourrai jamais te quitter je t'en} \\ \underline{fais la promesse mon bébé. Tu m'as tellement épaulé quand il le fallait. Tu m'as sorti de la} \\ \underline{galère je ne fais que rendre la monnaie.}
\begin{itemize}
\item Si jamais c'était possible, j'enlèverai toutes les épines de ta rose juste pour te serrer dans mes bras
\item J'ai beau me dire que je ne suis pas de taille, mais je fais semblant devant toi comme si c'était normal.
\item It is for you babe.
\end{itemize}

\rule{16cm}{1pt} \\ \\
{\Large \underline{\textbf{TP 3 :}} Section et structurations} 
\section{Ma famille}
\subsection{Papa}
\subsection{Maman}

\newpage
{\Large \underline{\textbf{TP 4 :}} Insertion d'équations} \\ \\
Comme exemple d'équations simples, on peut avoir : $x^2+2x=x-12$ \\ \\
Une fraction, on peut avoir : $\frac{5}{9}$ ou $\frac{x+1}{2}$ \\ \\
Une équation mathématique à racine carrée : $\sqrt{x^2+2x+3}$ \\ \\
\rule{16cm}{1pt} \\ \\

{\Large \underline{\textbf{TP 5 :}} Tableaux }
\begin{table}[h]
\begin{tabular}{|p{2cm}|p{2.5cm}|}
\hline
\textbf{MANGAS} & \textbf{SÉRIES} \\
\hline  
SNK & LAST SHIP \\
\hline
Naruto & Suits \\
\hline
\end{tabular}
\centering \caption {Réalisé par Gon FREECSS}
\end{table}
\\ 
\textbf{TP 6 :} Insertion d'images
\begin{figure}[h]
	\centering
	\includegraphics[width=0.8\textwidth]{azer}
	\caption{My life}
\end{figure}
\begin{enumerate}
\item Donne 
\begin{enumerate}
\item jhr
\end{enumerate}
\end{enumerate}
\end{document}